% ============================================================
% Question Bank: ch09-05 Sample Spaces for Compound Events
% Topic Title: Sample Spaces for Compound Events
% CCSS: 7.SP.C.8a
% Questions: 27 (15 MC + 6 SA + 6 Graphical)
% ============================================================

% --- Multiple Choice Questions (15) ---

\begin{practiceQuestion}{9-5-q01}{mc}
\begin{questionText}
A coin is flipped and a spinner with $5$ equal sections (A, B, C, D, E) is spun. How many outcomes are in the sample space?
\end{questionText}
\choiceA{$5$}
\choiceB{$7$}
\choiceC{$10$}
\choiceD{$25$}
\correctAnswer{C}
\explanation{The total number of outcomes is the product of the individual counts: $2 \times 5 = 10$. Each of the $2$ coin outcomes can pair with each of the $5$ spinner sections.}
\end{practiceQuestion}

\begin{practiceQuestion}{9-5-q02}{mc}
\begin{questionText}
A pizza shop offers $4$ crust types and $3$ sauce options. If a customer picks one crust and one sauce, how many different combinations are possible?
\end{questionText}
\choiceA{$7$}
\choiceB{$12$}
\choiceC{$16$}
\choiceD{$24$}
\correctAnswer{B}
\explanation{Using the counting principle, the total combinations are $4 \times 3 = 12$. Each crust type pairs with each sauce option.}
\end{practiceQuestion}

\begin{practiceQuestion}{9-5-q03}{mc}
\begin{questionText}
Two standard number cubes are rolled. How many outcomes are in the sample space?
\end{questionText}
\choiceA{$6$}
\choiceB{$12$}
\choiceC{$24$}
\choiceD{$36$}
\correctAnswer{D}
\explanation{Each cube has $6$ faces. The total outcomes for rolling two cubes is $6 \times 6 = 36$.}
\end{practiceQuestion}

\begin{practiceQuestion}{9-5-q04}{mc}
\begin{questionText}
Which method is best for displaying ALL outcomes of flipping a coin and rolling a number cube?
\end{questionText}
\choiceA{A bar graph}
\choiceB{A circle graph}
\choiceC{A tree diagram}
\choiceD{A line graph}
\correctAnswer{C}
\explanation{A tree diagram systematically branches out all possibilities, making it the best tool for listing every outcome in a compound event's sample space.}
\end{practiceQuestion}

\begin{practiceQuestion}{9-5-q05}{mc}
\begin{questionText}
A shirt comes in $3$ sizes (S, M, L) and $6$ colors. How many different shirts are available?
\end{questionText}
\choiceA{$9$}
\choiceB{$18$}
\choiceC{$36$}
\choiceD{$12$}
\correctAnswer{B}
\explanation{By the counting principle, the number of combinations is $3 \times 6 = 18$.}
\end{practiceQuestion}

\begin{practiceQuestion}{9-5-q06}{mc}
\begin{questionText}
Three coins are flipped. Which list shows the complete sample space?
\end{questionText}
\choiceA{HHH, HHT, HTH, THH, TTT}
\choiceB{HHH, HHT, HTH, HTT, THH, THT, TTH, TTT}
\choiceC{HH, HT, TH, TT}
\choiceD{HHH, HHT, HTT, TTT}
\correctAnswer{B}
\explanation{Three coins produce $2^3 = 8$ outcomes. Choice B lists all $8$ systematically: HHH, HHT, HTH, HTT, THH, THT, TTH, TTT.}
\end{practiceQuestion}

\begin{practiceQuestion}{9-5-q07}{mc}
\begin{questionText}
A school lunch offers a choice of $2$ entrees (chicken or fish) and $4$ drinks (milk, juice, water, tea). Which outcome is NOT possible?
\end{questionText}
\choiceA{Chicken and milk}
\choiceB{Fish and juice}
\choiceC{Chicken and soda}
\choiceD{Fish and tea}
\correctAnswer{C}
\explanation{Soda is not one of the $4$ drink options (milk, juice, water, tea), so ``Chicken and soda'' is not in the sample space.}
\end{practiceQuestion}

\begin{practiceQuestion}{9-5-q08}{mc}
\begin{questionText}
A locker combination uses one letter from $\{A, B, C\}$ followed by one digit from $\{1, 2, 3, 4\}$. How many possible combinations are there?
\end{questionText}
\choiceA{$7$}
\choiceB{$12$}
\choiceC{$16$}
\choiceD{$24$}
\correctAnswer{B}
\explanation{Using the counting principle: $3$ letters $\times$ $4$ digits $= 12$ possible combinations.}
\end{practiceQuestion}

\begin{practiceQuestion}{9-5-q09}{mc}
\begin{questionText}
A number cube is rolled and a coin is flipped. Zara uses a table to organize the sample space. How many cells should the table have?
\end{questionText}
\choiceA{$6$}
\choiceB{$8$}
\choiceC{$12$}
\choiceD{$36$}
\correctAnswer{C}
\explanation{The table should have $6$ rows (one for each die face) and $2$ columns (heads, tails), giving $6 \times 2 = 12$ cells for the $12$ outcomes.}
\end{practiceQuestion}

\begin{practiceQuestion}{9-5-q10}{mc}
\begin{questionText}
A bag has $2$ marbles (red and blue). You pick one marble, replace it, and pick again. Using an organized list, how many outcomes are in the sample space?
\end{questionText}
\choiceA{$2$}
\choiceB{$3$}
\choiceC{$4$}
\choiceD{$6$}
\correctAnswer{C}
\explanation{With replacement, each draw has $2$ options. The outcomes are: RR, RB, BR, BB — $2 \times 2 = 4$ total.}
\end{practiceQuestion}

\begin{practiceQuestion}{9-5-q11}{mc}
\begin{questionText}
An ice cream shop offers vanilla, chocolate, and strawberry ice cream in a cup or a cone. Which is the total number of options for one scoop?
\end{questionText}
\choiceA{$3$}
\choiceB{$5$}
\choiceC{$6$}
\choiceD{$9$}
\correctAnswer{C}
\explanation{By the counting principle: $3$ flavors $\times$ $2$ containers $= 6$ options.}
\end{practiceQuestion}

\begin{practiceQuestion}{9-5-q12}{mc}
\begin{questionText}
A spinner has $4$ equal sections (1, 2, 3, 4). The spinner is spun twice. How many outcomes are in the sample space?
\end{questionText}
\choiceA{$4$}
\choiceB{$8$}
\choiceC{$12$}
\choiceD{$16$}
\correctAnswer{D}
\explanation{Each spin has $4$ outcomes. Spinning twice: $4 \times 4 = 16$ total outcomes in the sample space.}
\end{practiceQuestion}

\begin{practiceQuestion}{9-5-q13}{mc}
\begin{questionText}
A car dealership offers $5$ models, $4$ exterior colors, and $2$ interior colors. How many different car configurations are possible?
\end{questionText}
\choiceA{$11$}
\choiceB{$20$}
\choiceC{$40$}
\choiceD{$80$}
\correctAnswer{C}
\explanation{Using the counting principle: $5 \times 4 \times 2 = 40$ configurations.}
\end{practiceQuestion}

\begin{practiceQuestion}{9-5-q14}{mc}
\begin{questionText}
A game involves drawing a card labeled $1$, $2$, or $3$ and then flipping a coin. One student says there are $5$ outcomes; another says $6$. Who is correct?
\end{questionText}
\choiceA{The first student, because $3 + 2 = 5$.}
\choiceB{The second student, because $3 \times 2 = 6$.}
\choiceC{Neither — there are $9$ outcomes.}
\choiceD{Neither — there are $3$ outcomes.}
\correctAnswer{B}
\explanation{For compound events, you multiply (not add) the number of choices. $3$ cards $\times$ $2$ coin sides $= 6$ outcomes. The second student is correct.}
\end{practiceQuestion}

\begin{practiceQuestion}{9-5-q15}{mc}
\begin{questionText}
Five coins are flipped. How many outcomes are in the sample space?
\end{questionText}
\choiceA{$10$}
\choiceB{$16$}
\choiceC{$25$}
\choiceD{$32$}
\correctAnswer{D}
\explanation{Each coin has $2$ outcomes. For $5$ coins: $2^5 = 32$ total outcomes.}
\end{practiceQuestion}

% --- Short Answer Questions (6) ---

\begin{practiceQuestion}{9-5-q16}{sa}
\begin{questionText}
A cafeteria offers $3$ main dishes (pasta, salad, sandwich) and $5$ beverages. How many meal combinations (one main dish and one beverage) are possible?
\end{questionText}
\correctAnswer{$15$}
\explanation{By the counting principle: $3 \times 5 = 15$ combinations.}
\end{practiceQuestion}

\begin{practiceQuestion}{9-5-q17}{sa}
\begin{questionText}
A bag has $4$ marbles (red, blue, green, yellow). You draw one marble, do NOT replace it, and draw another. How many outcomes are in the sample space?
\end{questionText}
\correctAnswer{$12$}
\explanation{Without replacement, the first draw has $4$ choices and the second has $3$: $4 \times 3 = 12$ outcomes.}
\end{practiceQuestion}

\begin{practiceQuestion}{9-5-q18}{sa}
\begin{questionText}
A security code consists of one digit ($0$--$9$) followed by one letter (A--Z). How many possible codes are there?
\end{questionText}
\correctAnswer{$260$}
\explanation{There are $10$ digit choices and $26$ letter choices. By the counting principle: $10 \times 26 = 260$.}
\end{practiceQuestion}

\begin{practiceQuestion}{9-5-q19}{sa}
\begin{questionText}
A coin is flipped $4$ times. How many outcomes are in the sample space?
\end{questionText}
\correctAnswer{$16$}
\explanation{Each flip has $2$ outcomes. For $4$ flips: $2^4 = 16$ outcomes.}
\end{practiceQuestion}

\begin{practiceQuestion}{9-5-q20}{sa}
\begin{questionText}
A quiz has $3$ true/false questions. List all the outcomes in the sample space. Use T for true and F for false.
\end{questionText}
\correctAnswer{TTT, TTF, TFT, TFF, FTT, FTF, FFT, FFF}
\explanation{Each question has $2$ answers, so $2^3 = 8$ outcomes. They are: TTT, TTF, TFT, TFF, FTT, FTF, FFT, FFF.}
\end{practiceQuestion}

\begin{practiceQuestion}{9-5-q21}{sa}
\begin{questionText}
A phone cover comes in $6$ designs and $3$ materials. If the company adds $2$ more designs, how many total options will there be?
\end{questionText}
\correctAnswer{$24$}
\explanation{After adding $2$ designs: $6 + 2 = 8$ designs. Total options: $8 \times 3 = 24$.}
\end{practiceQuestion}

% --- Graphical Multiple Choice Questions (3) ---

\begin{practiceQuestion}{9-5-q22}{gmc}
\begin{questionText}
Look at the tree diagram below showing the outcomes of picking a snack and a drink.

\begin{center}
\begin{tikzpicture}[x=1cm, y=1cm, >=Latex, font=\small]
  \node at (0,4.2) {Start};
  \node at (-3,2.7) {Apple};
  \node at (0,2.7) {Granola};
  \node at (3,2.7) {Crackers};
  \node at (-4,0.9) {Water};
  \node at (-2,0.9) {Juice};
  \node at (-1,0.9) {Water};
  \node at (1,0.9) {Juice};
  \node at (2,0.9) {Water};
  \node at (4,0.9) {Juice};

  \draw[thick,->] (0,3.95) -- (-3,3.02);
  \draw[thick,->] (0,3.95) -- (0,3.02);
  \draw[thick,->] (0,3.95) -- (3,3.02);

  \draw[thick,->] (-3,2.42) -- (-4,1.18);
  \draw[thick,->] (-3,2.42) -- (-2,1.18);
  \draw[thick,->] (0,2.42) -- (-1,1.18);
  \draw[thick,->] (0,2.42) -- (1,1.18);
  \draw[thick,->] (3,2.42) -- (2,1.18);
  \draw[thick,->] (3,2.42) -- (4,1.18);
\end{tikzpicture}
\end{center}

How many outcomes are in the sample space?
\end{questionText}
\choiceA{$3$}
\choiceB{$5$}
\choiceC{$6$}
\choiceD{$9$}
\correctAnswer{C}
\explanation{The tree diagram shows $3$ snack choices each paired with $2$ drink choices. Counting the branch endpoints gives $3 \times 2 = 6$ outcomes.}
\end{practiceQuestion}

\begin{practiceQuestion}{9-5-q23}{gmc}
\begin{questionText}
The table below shows the sample space for rolling a number cube and flipping a coin. Some outcomes are filled in.

\begin{center}
\begin{tabular}{c|cc}
 & \textbf{Heads} & \textbf{Tails} \\
\hline
\textbf{1} & (1, H) & (1, T) \\
\textbf{2} & (2, H) & (2, T) \\
\textbf{3} & (3, H) & ? \\
\textbf{4} & (4, H) & (4, T) \\
\textbf{5} & ? & (5, T) \\
\textbf{6} & (6, H) & (6, T) \\
\end{tabular}
\end{center}

How many total outcomes are in this sample space?
\end{questionText}
\choiceA{$6$}
\choiceB{$8$}
\choiceC{$12$}
\choiceD{$36$}
\correctAnswer{C}
\explanation{The table has $6$ rows (die faces) and $2$ columns (coin sides). The total number of outcomes is $6 \times 2 = 12$.}
\end{practiceQuestion}

\begin{practiceQuestion}{9-5-q24}{gmc}
\begin{questionText}
A tree diagram shows outcomes for choosing a pen color (Black or Blue) and a pen type (Gel, Ballpoint, or Felt-tip).

\begin{center}
\begin{tikzpicture}[x=1cm, y=1cm, >=Latex, font=\small]
  \node at (0,4.4) {Start};
  \node at (-2.8,2.8) {Black};
  \node at (2.8,2.8) {Blue};
  \node at (-4.6,1.0) {Gel};
  \node at (-2.8,1.0) {Ballpoint};
  \node at (-1.0,1.0) {Felt-tip};
  \node at (1.0,1.0) {Gel};
  \node at (2.8,1.0) {Ballpoint};
  \node at (4.6,1.0) {Felt-tip};

  \draw[thick,->] (0,4.12) -- (-2.8,3.08);
  \draw[thick,->] (0,4.12) -- (2.8,3.08);

  \draw[thick,->] (-2.8,2.52) -- (-4.6,1.28);
  \draw[thick,->] (-2.8,2.52) -- (-2.8,1.28);
  \draw[thick,->] (-2.8,2.52) -- (-1.0,1.28);
  \draw[thick,->] (2.8,2.52) -- (1.0,1.28);
  \draw[thick,->] (2.8,2.52) -- (2.8,1.28);
  \draw[thick,->] (2.8,2.52) -- (4.6,1.28);
\end{tikzpicture}
\end{center}

Which outcome is NOT shown in this tree diagram?
\end{questionText}
\choiceA{Black Gel}
\choiceB{Blue Ballpoint}
\choiceC{Red Felt-tip}
\choiceD{Blue Gel}
\correctAnswer{C}
\explanation{The tree diagram only includes Black and Blue pen colors. Red is not one of the options, so ``Red Felt-tip'' is not in the sample space.}
\end{practiceQuestion}

% --- Graphical Short Answer Questions (3) ---

\begin{practiceQuestion}{9-5-q25}{gsa}
\begin{questionText}
Complete the table to show the sample space for spinning a spinner with sections $\{X, Y, Z\}$ and rolling a die with faces $\{1, 2\}$.

\begin{center}
\begin{tabular}{c|cc}
 & \textbf{1} & \textbf{2} \\
\hline
\textbf{X} & (X, 1) & ? \\
\textbf{Y} & ? & (Y, 2) \\
\textbf{Z} & (Z, 1) & ? \\
\end{tabular}
\end{center}

How many total outcomes are in this sample space?
\end{questionText}
\correctAnswer{$6$}
\explanation{The table has $3$ rows and $2$ columns, giving $3 \times 2 = 6$ outcomes: (X,1), (X,2), (Y,1), (Y,2), (Z,1), (Z,2).}
\end{practiceQuestion}

\begin{practiceQuestion}{9-5-q26}{gsa}
\begin{questionText}
Draw the sample space as an organized list for the compound event: flip a coin and pick a card from $\{1, 2, 3\}$.

\begin{center}
\begin{tikzpicture}[x=1cm, y=1cm, >=Latex, font=\small]
  \node at (0,4.4) {Start};
  \node at (-2.3,2.8) {H};
  \node at (2.3,2.8) {T};
  \node at (-3.9,0.9) {H1};
  \node at (-2.3,0.9) {H2};
  \node at (-0.7,0.9) {H3};
  \node at (0.7,0.9) {T1};
  \node at (2.3,0.9) {T2};
  \node at (3.9,0.9) {T3};

  \draw[thick,->] (0,4.12) -- (-2.3,3.06);
  \draw[thick,->] (0,4.12) -- (2.3,3.06);

  \draw[thick,->] (-2.3,2.48) -- (-3.9,1.18);
  \draw[thick,->] (-2.3,2.48) -- (-2.3,1.18);
  \draw[thick,->] (-2.3,2.48) -- (-0.7,1.18);
  \draw[thick,->] (2.3,2.48) -- (0.7,1.18);
  \draw[thick,->] (2.3,2.48) -- (2.3,1.18);
  \draw[thick,->] (2.3,2.48) -- (3.9,1.18);

  \node[font=\scriptsize] at (-3.15,1.95) {1};
  \node[font=\scriptsize] at (-2.3,1.9) {2};
  \node[font=\scriptsize] at (-1.45,1.95) {3};
  \node[font=\scriptsize] at (1.45,1.95) {1};
  \node[font=\scriptsize] at (2.3,1.9) {2};
  \node[font=\scriptsize] at (3.15,1.95) {3};
\end{tikzpicture}
\end{center}

How many outcomes have tails and an even-numbered card?
\end{questionText}
\correctAnswer{$1$}
\explanation{From the tree diagram, the outcomes with tails are T1, T2, T3. The only even number is $2$, so the only outcome with tails and an even card is T2. That gives $1$ outcome.}
\end{practiceQuestion}

\begin{practiceQuestion}{9-5-q27}{gsa}
\begin{questionText}
The sample space table below shows the sums when two number cubes are rolled.

\begin{center}
\begin{tabular}{c|cccccc}
$+$ & \textbf{1} & \textbf{2} & \textbf{3} & \textbf{4} & \textbf{5} & \textbf{6} \\
\hline
\textbf{1} & $2$ & $3$ & $4$ & $5$ & $6$ & $7$ \\
\textbf{2} & $3$ & $4$ & $5$ & $6$ & $7$ & $8$ \\
\textbf{3} & $4$ & $5$ & $6$ & $7$ & $8$ & $9$ \\
\textbf{4} & $5$ & $6$ & $7$ & $8$ & $9$ & $10$ \\
\textbf{5} & $6$ & $7$ & $8$ & $9$ & $10$ & $11$ \\
\textbf{6} & $7$ & $8$ & $9$ & $10$ & $11$ & $12$ \\
\end{tabular}
\end{center}

How many outcomes produce a sum of $9$ or greater?
\end{questionText}
\correctAnswer{$10$}
\explanation{Count the cells with sums $\geq 9$: sum $9$ appears $4$ times (3+6, 4+5, 5+4, 6+3), sum $10$ appears $3$ times (4+6, 5+5, 6+4), sum $11$ appears $2$ times (5+6, 6+5), sum $12$ appears $1$ time (6+6). Total: $4 + 3 + 2 + 1 = 10$.}
\end{practiceQuestion}
